% Book 4, Year 2, chapter 30 -- Database design: dependencies, normal forms and transactions.
% Provenance (never printed): sources/book4/BRIEFS.md, chapter 30.
\chapter{Database Design: Dependencies, Normal Forms and Transactions}\label{ch:b2:database-design}

For years the library kept its loans in one spreadsheet: one line per loan, with the book's title and author, the
member's name and street, everything a librarian might want to read on the line. Built as a table from the town's
database, it looks like this, and it shows what happens when member 55 moves house:

\omcode{code/bachelor-2/30-database-design/listings/wide_table.sql}{1}{10}{The wide table, and a move recorded on every
line but the last}\label{lst:b2:database-design:hook}

\omoutput{code/bachelor-2/30-database-design/outputs/wide_table.txt}{Eleven lines changed, one forgotten: the member now
lives in two streets at once.}\label{out:b2:database-design:hook}

The listing loads \code{wide_build.sql}, which builds the table by joining loans, books, members, people and streets.
Nothing in the table says that a member has one street: the fact is written twelve times, and the twelve copies can
disagree. This chapter makes the intuition precise. Dependencies say which facts determine which; normal forms say
when a table stores each fact once; \omterm{def:b2:database-design:decomposition}{decompositions} split a table without losing anything. The second half turns to the
other way a database goes wrong, two programs changing it at the same moment, and ends with the \omterm{def:b2:database-design:index}{index} that makes a
search fast.

\begin{recall}
The Year 1 volume: relations, attributes, keys, candidate and primary keys, foreign keys, relational algebra and its
natural join, SQL queries, entity--relationship modelling. The grade 12 volume: the anomalies of a badly designed table,
seen on examples.
\end{recall}

\section{Functional dependencies}\label{sec:b2:database-design:fd}

\begin{definition}[Functional dependency, trivial dependency]\label{def:b2:database-design:fd}
Let $R$ be a set of attributes and $X, Y \subseteq R$. A table over $R$ \emph{satisfies} the \emph{functional
dependency}\index{functional dependency} $X \to Y$ if any two rows that agree on every attribute of $X$ also agree on every
attribute of $Y$. The dependency is \emph{trivial}\index{trivial dependency} if $Y \subseteq X$; it is then satisfied by
every table.
\end{definition}

A dependency is a rule about the world the table describes, not an observation about today's rows. The library's rules
are that a loan determines its book, its member and its dates; a book determines its title, author, genre and pages; a
member determines the name and the street; a street determines its postcode. The 400 lines satisfy other dependencies
by accident: no two books have the same title yet, so \emph{title} $\to$ \emph{book} holds in the data, but nothing
forbids two books called the same (\cref{exo:b2:database-design:9}).

\begin{omfigure}
\begin{tikzpicture}[a/.style={draw, rounded corners=2pt, font=\scriptsize\ttfamily, inner sep=2.5pt,
  text height=1.4ex, text depth=0.4ex}, d/.style={->, >=Stealth, thick}]
\node[a] (loan) at (0, 0) {loan};
\foreach \n [remember=\n as \p (initially loan)] in {out, due, ret, book, title, author, genre, pages, member,
  first, last, street, postcode} \node[a, right=1.6mm of \p] (\n) {\n};
\coordinate (mIn) at ([xshift=-1.5mm]member.north); \coordinate (mOut) at ([xshift=1.5mm]member.north);
\coordinate (blue) at ($(loan.north) + (0, 0.45)$); \coordinate (orange) at ($(mOut) + (0, 0.9)$);
\coordinate (red) at ($(book.south) + (0, -0.45)$);
\draw[omDef, thick] (loan.north) -- (blue) -- (blue -| mIn);
\foreach \t in {out, due, ret, book} \draw[d, omDef] (blue -| \t.north) -- (\t.north);
\draw[d, omDef] (blue -| mIn) -- (mIn);
\draw[omProp, thick] (mOut) -- (orange) -- (orange -| street.north);
\foreach \t in {first, last, street} \draw[d, omProp] (orange -| \t.north) -- (\t.north);
\draw[omThm, thick] (book.south) -- (red) -- (red -| pages.south);
\foreach \t in {title, author, genre, pages} \draw[d, omThm] (red -| \t.south) -- (\t.south);
\draw[d, omMeth] (street.south) -- ++(0, -0.45) -| (postcode.south);
\end{tikzpicture}
\par\smallskip{\small The library's rules over the wide table's fourteen attributes. Blue: a loan determines its dates,
its book and its member. Red: a book determines its title, author, genre and pages. Orange: a member determines the name
and the street. Green: a street determines its postcode. Only the blue arrows start from a key.}
\end{omfigure}

From some dependencies others follow: if a member determines the street and a street the postcode, a member determines
the postcode.

\begin{definition}[Implication, Armstrong's axioms]\label{def:b2:database-design:armstrong}
A set $F$ of dependencies over $R$ \emph{implies}\index{implication of dependencies} $X \to Y$, written $F \models X \to
Y$, if every table over $R$ that satisfies all of $F$ satisfies $X \to Y$. \emph{Armstrong's
axioms}\index{Armstrong's axioms} are three rules to derive dependencies, written $F \vdash X \to Y$:
\emph{reflexivity}, $X \to Y$ whenever $Y \subseteq X$; \emph{augmentation}, from $X \to Y$ derive $XZ \to YZ$;
\emph{transitivity}, from $X \to Y$ and $Y \to Z$ derive $X \to Z$. (We write $XZ$ for $X \cup Z$.)
\end{definition}

The \emph{union} rule, from $X \to Y$ and $X \to Z$ derive $X \to YZ$, and the \emph{decomposition} rule, from $X \to YZ$
derive $X \to Y$, follow from the three (\cref{exo:b2:database-design:8}).

\begin{definition}[Attribute closure]\label{def:b2:database-design:closure}
The \emph{attribute closure}\index{attribute closure} of $X$ under $F$ is $X^+ = \{A \in R : F \models X \to A\}$, the
set of attributes that $X$ determines.
\end{definition}

\begin{theorem}[Armstrong's axioms are sound and complete]\label{thm:b2:database-design:armstrong}
For every $F$, $X$ and $Y$: $F \vdash X \to Y$ if and only if $F \models X \to Y$.
\end{theorem}

\begin{proof}
\emph{Soundness}: each rule turns satisfied dependencies into a satisfied one. Reflexivity: two rows that agree on $X$
agree on its subset $Y$. Augmentation: two rows that agree on $XZ$ agree on $X$, hence on $Y$, and on $Z$. Transitivity:
two rows that agree on $X$ agree on $Y$, hence on $Z$. By induction on the length of a derivation, every derived
dependency holds in every table satisfying $F$.

\emph{Completeness}: suppose $F \models X \to Y$, and let $D = \{A : F \vdash X \to A\}$; by reflexivity $X \subseteq D$.
Build a table of two rows over $R$: both rows are $0$ on $D$; outside $D$, the first is $0$ and the second $1$. It
satisfies every $V \to W$ of $F$: if the rows agree on $V$, then $V \subseteq D$, so $F \vdash X \to V$ by the union rule,
then $F \vdash X \to W$ by transitivity, so $W \subseteq D$ by decomposition and the rows agree on $W$. The table
therefore satisfies $X \to Y$; its rows agree on $X \subseteq D$, so they agree on $Y$, and $Y \subseteq D$. By the union
rule, $F \vdash X \to Y$.
\end{proof}

So the closure can be computed by derivation instead of by imagining every table, and there is a simple algorithm:

\omcode{code/bachelor-2/30-database-design/listings/fd.py}{15}{24}{The closure of a set of attributes: apply every
dependency whose left side is already in, until nothing changes}\label{lst:b2:database-design:closure}

\begin{proposition}[The closure algorithm]\label{prop:b2:database-design:closure-algo}
The function \code{closure} returns $X^+$, after at most $|F| + 1$ passes over $F$, so in time $O(|F|^2 \cdot |R|)$.
\end{proposition}

\begin{proof}
\Cref{exo:b2:database-design:11}: every attribute added is derivable, and when the loop stops the set is closed under
$F$, which forces it to contain $X^+$. Each pass but the last adds the right side of at least one dependency not used
before.
\end{proof}

\begin{corollary}[A test for keys]\label{cor:b2:database-design:key-test}
$X$ is a superkey of $R$ under $F$ exactly when $X^+ = R$; it is a candidate key when moreover $(X \setminus \{A\})^+ \neq R$
for every $A \in X$.
\end{corollary}

\begin{proof}
$X$ is a superkey when it determines every attribute, that is $F \models X \to R$, that is $R \subseteq X^+$. A candidate
key is a superkey with no proper subset that is one; since a superset of a superkey is a superkey, it suffices to remove
one attribute at a time.
\end{proof}

\begin{definition}[Minimal cover]\label{def:b2:database-design:cover}
A \emph{minimal cover}\index{minimal cover} of $F$ is a set $G$ of dependencies equivalent to $F$ (each \omterm{def:b2:database-design:armstrong}{implies} every
dependency of the other) in which every right side is a single attribute, no attribute can be removed from a left side,
and no dependency can be removed, without losing the equivalence.
\end{definition}

\omcode{code/bachelor-2/30-database-design/listings/fd.py}{31}{57}{Candidate keys by increasing size, and a \omterm{def:b2:database-design:cover}{minimal
cover}}\label{lst:b2:database-design:keys}

\omcode{code/bachelor-2/30-database-design/listings/fd_demo.py}{5}{8}{The wide table's attributes and the library's
rules}\label{lst:b2:database-design:rules}

\omoutput{code/bachelor-2/30-database-design/outputs/fd_demo.txt}{Closures, the key, the cover, and the BCNF tables of
\cref{sec:b2:database-design:decomposition}.}\label{out:b2:database-design:fd-demo}

The key search tries the subsets in increasing size, skipping supersets of keys already found: exponential in general,
fast here, where one attribute suffices. The wide table's only candidate key is the loan.

\section{Normal forms}\label{sec:b2:database-design:normal-forms}

\begin{definition}[Update, insertion and deletion anomaly]\label{def:b2:database-design:anomaly}
A table has an \emph{update anomaly}\index{update anomaly} when one fact is stored on several rows, so that changing it
means changing them all; an \emph{insertion anomaly}\index{insertion anomaly} when a fact cannot be recorded without an
unrelated one, as a member without a loan; a \emph{deletion anomaly}\index{deletion anomaly} when removing a row removes
an unrelated fact with it, as the last loan of a book removes the book.
\end{definition}

All three come from a dependency whose left side is not a key: \emph{book} $\to$ \emph{title} makes the title a fact
about the book, stored once per loan of that book.

\begin{definition}[Third normal form, Boyce--Codd normal form]\label{def:b2:database-design:normal-forms}
Let $F$ hold on $R$. $R$ is in \emph{Boyce--Codd normal form}\index{Boyce--Codd normal form} (BCNF) if for every
\omterm{def:b2:database-design:fd}{non-trivial} $X \to A$ implied by $F$ with $X \cup \{A\} \subseteq R$, $X$ is a superkey of $R$. It is in \emph{third
normal form}\index{third normal form} (3NF) if for every such dependency, $X$ is a superkey or $A$ belongs to some
candidate key of $R$.
\end{definition}

BCNF says that every fact the table records is a fact about a key; 3NF tolerates one exception, attributes that are
part of a key. The wide table is in neither: \emph{book}, \emph{member} and \emph{street} determine other attributes and
are not keys.

\section{Decomposition}\label{sec:b2:database-design:decomposition}

\begin{definition}[Decomposition, lossless join, dependency preservation]\label{def:b2:database-design:decomposition}
A \emph{decomposition}\index{decomposition} of $R$ is a list of attribute sets $R_1, \dots, R_m$ whose union is $R$; a table
$T$ over $R$ is stored as its projections $\pi_{R_i}(T)$. The decomposition is
\emph{lossless}\index{lossless-join decomposition} for $F$ if $T = \pi_{R_1}(T) \bowtie \dots \bowtie \pi_{R_m}(T)$ for
every $T$ satisfying $F$; it \emph{preserves the dependencies}\index{dependency preservation} if the dependencies of $F$
that hold inside the parts imply all of $F$.
\end{definition}

The join of the projections always contains $T$; a lossy \omterm{def:b2:database-design:decomposition}{decomposition} adds rows that were never there. Preservation
matters for checking: a dependency that holds inside a part is enforced by that part's key, while one that spans two
parts can be checked only by a join.

\begin{theorem}[Heath's theorem]\label{thm:b2:database-design:heath}
Let $R = X \cup Y \cup Z$ with $X, Y, Z$ disjoint. If a table over $R$ satisfies $X \to Y$, it is the join of its
projections on $XY$ and $XZ$.
\end{theorem}

\begin{proof}
\Cref{exo:b2:database-design:10}.
\end{proof}

\needspace{18\baselineskip}
\omcode{code/bachelor-2/30-database-design/listings/fd.py}{67}{83}{A BCNF violation, and the BCNF \omterm{def:b2:database-design:decomposition}{decomposition} by
repeated splits}\label{lst:b2:database-design:bcnf}

\begin{theorem}[BCNF decomposition]\label{thm:b2:database-design:bcnf-lossless}
The function \code{bcnf} stops, returns parts that are all in BCNF, and the \omterm{def:b2:database-design:decomposition}{decomposition} is \omterm{def:b2:database-design:decomposition}{lossless}.
\end{theorem}

\begin{proof}
A split of $R$ on a violation $X \to Y$, with $Y = X^+ \cap R \setminus X$ non-empty and $X^+ \cap R \neq R$, produces $XY$
and $R \setminus Y$, both strictly smaller than $R$: the recursion stops. A part is returned only when it has no violation,
so it is in BCNF. Each split is \omterm{def:b2:database-design:decomposition}{lossless} by Heath's theorem with $Z = R \setminus XY$, since $X \to Y$ holds. By
induction on the recursion, if $R \setminus Y$ is replaced by a \omterm{def:b2:database-design:decomposition}{lossless} \omterm{def:b2:database-design:decomposition}{decomposition} of it, the whole stays \omterm{def:b2:database-design:decomposition}{lossless},
because a table satisfying $F$ projects onto $R \setminus Y$ as a table that satisfies the dependencies of $F$ that hold
there, and the natural join is associative.
\end{proof}

\begin{omfigure}
\begin{forest}
for tree={font=\scriptsize, draw, rounded corners, align=center, l sep=14pt, s sep=8pt, inner sep=3pt},
[{wide: 14 attributes, key loan}
  [{book, title, author,\\genre, pages}, edge label={node[midway, left, font=\tiny] {on book}}]
  [{loan, dates, book,\\member, first, last,\\street, postcode}
    [{member, first, last,\\street, postcode}, edge label={node[midway, left, font=\tiny] {on member}}
      [{street, postcode}, edge label={node[midway, left, font=\tiny] {on street}}]
      [{member, first,\\last, street}]]
    [{loan, out, due, ret,\\book, member}]]]
\end{forest}
\par\smallskip{\small The BCNF \omterm{def:b2:database-design:decomposition}{decomposition} of \cref{out:b2:database-design:fd-demo}. Each split takes a violating
dependency $X \to Y$ and makes one table of $X$ with everything it determines, and another of the rest with $X$ kept as
the link. The four leaves are the BCNF tables.}
\end{omfigure}

The SQL counterpart builds the four tables from the wide one and joins them back:

\omcode{code/bachelor-2/30-database-design/listings/decomposed.sql}{1}{23}{The four BCNF tables, and the join that should
give back the wide table}\label{lst:b2:database-design:decomposed}

\omoutput{code/bachelor-2/30-database-design/outputs/decomposed.txt}{No row lost, none invented, and each key unique in
the data.}\label{out:b2:database-design:decomposed}

Only 23 streets appear because no member lives on the 24th. BCNF cannot always keep the
dependencies inside the parts.

\begin{example}[No BCNF decomposition preserves the dependencies]\label{ex:b2:database-design:no-preserving-bcnf}
An address table over \emph{street}, \emph{town} and \emph{postcode} with the rules \emph{street town} $\to$
\emph{postcode} and \emph{postcode} $\to$ \emph{town}. The keys are \{street, town\} and \{street, postcode\}, so
\emph{postcode} $\to$ \emph{town} violates BCNF; every attribute belongs to a key, so the table is in 3NF. Any BCNF
\omterm{def:b2:database-design:decomposition}{decomposition} must separate the three attributes, and \emph{street town} $\to$ \emph{postcode} then lies in no part.
\end{example}

\omoutput{code/bachelor-2/30-database-design/outputs/no_preserving.txt}{The address table: BCNF loses a dependency, 3NF
keeps the table whole.}\label{out:b2:database-design:no-preserving}

\omcode{code/bachelor-2/30-database-design/listings/fd.py}{97}{106}{3NF synthesis from a \omterm{def:b2:database-design:cover}{minimal
cover}}\label{lst:b2:database-design:synthesis}

\begin{theorem}[3NF synthesis]\label{thm:b2:database-design:3nf-synthesis}
The \omterm{def:b2:database-design:decomposition}{decomposition} of \cref{lst:b2:database-design:synthesis}, one table for each left side of a \omterm{def:b2:database-design:cover}{minimal cover} with the
attributes it determines in the cover, plus a candidate key if no table contains one, \omterm{def:b2:database-design:decomposition}{preserves the dependencies}, is
\omterm{def:b2:database-design:decomposition}{lossless}, and has every table in 3NF.
\end{theorem}

\begin{proof}
Every dependency of the \omterm{def:b2:database-design:cover}{minimal cover} lies inside the table of its left side, and the cover is equivalent to $F$: the
dependencies are preserved. For losslessness, let $K$ be a table containing a key and run the chase of
\cref{exo:b2:database-design:5}: the row of $K$ starts with the symbol $a$ on $K$, and the cover's dependencies, applied in
the order in which the closure algorithm adds attributes to $K^+ = R$, give it $a$ on each attribute in turn, since each
such dependency lies inside a table whose own row has $a$ everywhere. A row with $a$ everywhere means the join is
\omterm{def:b2:database-design:decomposition}{lossless}. That each table is in 3NF uses the minimality of the cover; we admit it.
\end{proof}

So 3NF can always be reached without losing a dependency, and BCNF without losing information but perhaps a dependency.
The synthesis often gives BCNF tables anyway, as for the library: its four 3NF tables are the four BCNF tables.

One more kind of redundancy escapes functional dependencies. If a book has several authors and several subject
keywords, unrelated to each other, a table (\emph{book}, \emph{author}, \emph{keyword}) must hold every pair: two authors
and three keywords make six rows, with each author written three times.

\begin{definition}[Multivalued dependency, fourth normal form]\label{def:b2:database-design:mvd}
A table over $R$ satisfies the \emph{multivalued dependency}\index{multivalued dependency} $X \twoheadrightarrow Y$ if,
with $Z = R \setminus XY$, whenever two rows agree on $X$, the row that takes $X$ and $Y$ from the first and $Z$ from the
second is also in the table. $R$ is in \emph{fourth normal form}\index{fourth normal form} if for every \omterm{def:b2:database-design:fd}{non-trivial}
$X \twoheadrightarrow Y$ that holds, $X$ is a superkey.
\end{definition}

Here \emph{book} $\twoheadrightarrow$ \emph{author}: every author of a book appears with every keyword of it. The table
is in BCNF, its only key being all three attributes, but not in 4NF, and splitting it into (\emph{book}, \emph{author})
and (\emph{book}, \emph{keyword}) is \omterm{def:b2:database-design:decomposition}{lossless} by the same argument as Heath's theorem. A \omterm{def:b2:database-design:fd}{functional dependency} is a
special multivalued one; the chase of \cref{exo:b2:database-design:5} extends to both, and we leave it at that.

\begin{history}
% ledger: b2:codd-1970
% ledger: b2:codd-3nf-1971
% ledger: b2:bcnf-1974
Edgar Codd published the relational model in 1970 and defined \omterm{def:b2:database-design:normal-forms}{third normal form} in 1971. Raymond Boyce and Codd
formally developed the normal form that bears their names in 1974; Ian Heath had described it in 1971.
\end{history}

\section{Transactions}\label{sec:b2:database-design:transactions}

Design removes the anomalies of one table. The other dangers come from time: a change made of several statements that
stops in the middle, or two programs that change the same row at once.

\begin{definition}[Transaction, commit, rollback]\label{def:b2:database-design:transaction}
A \emph{transaction}\index{transaction} is a sequence of reads and writes that the database treats as one unit. It ends
with a \emph{commit}\index{commit}, which makes all its writes permanent and visible to others, or a
\emph{rollback}\index{rollback}, which cancels all of them.
\end{definition}

\begin{definition}[Atomicity, consistency, isolation, durability]\label{def:b2:database-design:acid}
A database system guarantees \emph{atomicity}\index{atomicity} if a \omterm{def:b2:database-design:transaction}{transaction}'s writes take effect all or none;
\emph{consistency}\index{consistency (transactions)} if a \omterm{def:b2:database-design:transaction}{transaction} that starts from a state satisfying the declared
constraints and \omterm{def:b2:database-design:transaction}{commits} leaves one; \emph{isolation}\index{isolation} if concurrent \omterm{def:b2:database-design:transaction}{transactions} do not see each other's
partial work, to a degree fixed by the \omterm{def:b2:database-design:isolation-level}{isolation level}; \emph{durability}\index{durability} if a committed \omterm{def:b2:database-design:transaction}{transaction}
survives a crash.
\end{definition}

In SQLite, \code{BEGIN} starts a \omterm{def:b2:database-design:transaction}{transaction}; without it, each statement is a \omterm{def:b2:database-design:transaction}{transaction} of its own. The book's SQL
runner \omterm{def:b2:database-design:transaction}{commits} after every statement, as an interactive tool does, so the \omterm{def:b2:database-design:transaction}{transactions} of this section are written in
Python, which controls them explicitly:

\omcode{code/bachelor-2/30-database-design/listings/rollback.py}{16}{25}{A late return recorded in two tables as one
\omterm{def:b2:database-design:transaction}{transaction}}\label{lst:b2:database-design:rollback}

\omoutput{code/bachelor-2/30-database-design/outputs/rollback.txt}{The fine of 0 breaks a constraint after the loan was
already updated: the \omterm{def:b2:database-design:transaction}{rollback} undoes the update too.}\label{out:b2:database-design:rollback}

Without the \omterm{def:b2:database-design:transaction}{transaction}, the first attempt would have left a loan returned with no fine recorded: the two tables would
disagree, the very thing the design was meant to prevent.

\begin{history}
% ledger: b2:haerder-reuter-1983
Theo Härder and Andreas Reuter coined the acronym ACID for these four properties in 1983, in a survey of \omterm{def:b2:database-design:transaction}{transaction}
recovery.
\end{history}

\section{Isolation}\label{sec:b2:database-design:isolation}

\begin{definition}[Schedule]\label{def:b2:database-design:schedule}
A \emph{schedule}\index{schedule (of transactions)} of a set of \omterm{def:b2:database-design:transaction}{transactions} is the order in time of all their reads,
writes, \omterm{def:b2:database-design:transaction}{commits} and \omterm{def:b2:database-design:transaction}{rollbacks}, each \omterm{def:b2:database-design:transaction}{transaction}'s own operations in their order. It is \emph{serial} if the \omterm{def:b2:database-design:transaction}{transactions}
do not overlap.
\end{definition}

\begin{definition}[The anomalies of isolation]\label{def:b2:database-design:anomalies}
In a \omterm{def:b2:database-design:schedule}{schedule} of two \omterm{def:b2:database-design:transaction}{transactions} $T_1$ and $T_2$: a \emph{dirty read}\index{dirty read} is $T_2$ reading a value written
by $T_1$ that $T_1$ has not committed; a \emph{non-repeatable read}\index{non-repeatable read} is $T_1$ reading the same
row twice and getting two values, because $T_2$ changed it and committed in between; a \emph{phantom
read}\index{phantom read} is $T_1$ running the same query with a condition twice and getting different sets of rows,
because $T_2$ inserted or deleted a matching row; a \emph{lost update}\index{lost update} is $T_1$ and $T_2$ both reading a
value and both writing a new value computed from it, so that one write erases the other; a \emph{write
skew}\index{write skew} is $T_1$ and $T_2$ each reading the same rows, checking a rule over them, and each writing a
different row, so that both checks pass and the result breaks the rule.
\end{definition}

\begin{definition}[Isolation level]\label{def:b2:database-design:isolation-level}
An \emph{isolation level}\index{isolation level} is a promise about which of these anomalies a database system lets
\omterm{def:b2:database-design:transaction}{transactions} see. The SQL standard names four: \emph{read uncommitted}, \emph{read committed}, \emph{repeatable read}
and \emph{serializable}, the last promising that every \omterm{def:b2:database-design:schedule}{schedule} has the same effect as some serial one.
\end{definition}

\begin{center}
% ledger: b2:sql-isolation-levels
\small
\begin{tabular}{lccc}
\hline
\omterm{def:b2:database-design:isolation-level}{isolation level} & \omterm{def:b2:database-design:anomalies}{dirty read} & \omterm{def:b2:database-design:anomalies}{non-repeatable read} & \omterm{def:b2:database-design:anomalies}{phantom read} \\
\hline
read uncommitted & possible & possible & possible \\
read committed & no & possible & possible \\
repeatable read & no & no & possible \\
serializable & no & no & no \\
\hline
\end{tabular}
\par\smallskip{\small The standard \omterm{def:b2:database-design:isolation-level}{isolation levels}, and the anomalies each allows. A system may give more than a level
promises: PostgreSQL's repeatable read shows no phantoms, and its read uncommitted behaves as read committed.}
\end{center}

SQLite takes the simplest road: its \omterm{def:b2:database-design:transaction}{transactions} are serializable, and it allows a single writer at a time. The one way
to see an uncommitted change is to share a cache between two connections and turn on \code{read_uncommitted}.

\begin{inpractice}
% ledger: b2:sqlite-serializable
% ledger: b2:sqlite-one-writer
SQLite's documentation states that its \omterm{def:b2:database-design:transaction}{transactions} are serializable except for connections that share a cache with
\code{read_uncommitted} on, and that one connection writes at a time: the others wait their turn.
\end{inpractice}

\omcode{code/bachelor-2/30-database-design/listings/dirty_read.py}{5}{27}{Two connections to one shared in-memory
database}\label{lst:b2:database-design:dirty-read}

\omoutput{code/bachelor-2/30-database-design/outputs/dirty_read.txt}{Refused, then a \omterm{def:b2:database-design:anomalies}{dirty read}, then the value the writer
never committed is gone.}\label{out:b2:database-design:dirty-read}

The \omterm{def:b2:database-design:anomalies}{lost update} needs no concurrency inside the database at all: it is enough that a program reads a value, computes,
and writes, while another does the same. Two lending desks, each with its own connection to the library's file:

\omcode{code/bachelor-2/30-database-design/listings/lost_update.py}{25}{59}{Read then write, unprotected, inside
\code{BEGIN IMMEDIATE}, then as one \code{UPDATE}}\label{lst:b2:database-design:lost-update}

\omoutput{code/bachelor-2/30-database-design/outputs/lost_update.txt}{Three copies, two lent: only the first way records
2 instead of 1.}\label{out:b2:database-design:lost-update}

\begin{omfigure}
\begin{tikzpicture}[x=1cm, y=0.62cm, s/.style={font=\small, anchor=west}]
\draw[thick, ->] (0, 0.4) -- (0, -5.4) node[below, font=\scriptsize] {time};
\node[font=\small\bfseries] at (2.6, 1) {desk A}; \node[font=\small\bfseries] at (7.6, 1) {desk B};
\node[font=\small\bfseries] at (11.6, 1) {shelf};
\node[s] at (1, 0) {reads \code{available}: 3}; \node[font=\small] at (11.6, 0) {3};
\node[s] at (6, -1) {reads \code{available}: 3}; \node[font=\small] at (11.6, -1) {3};
\node[s] at (1, -2) {writes $3 - 1 = 2$}; \node[font=\small] at (11.6, -2) {2};
\node[s] at (6, -3) {writes $3 - 1 = 2$}; \node[font=\small, text=omThm] at (11.6, -3) {2};
\draw[omThm, thick, ->] (8.6, -3.35) to[bend right=12] (4.2, -4.4) node[below, font=\small] {A's lending is lost};
\foreach \y in {0, -1, -2, -3} \draw[black!25] (0.9, \y - 0.5) -- (12.3, \y - 0.5);
\end{tikzpicture}
\par\smallskip{\small The \omterm{def:b2:database-design:anomalies}{lost update} of \cref{out:b2:database-design:lost-update}, one step per line. Each write is
correct for the value its desk read; the second erases the first. Inside \code{BEGIN IMMEDIATE}, desk B could not have
started before A committed, and it would have read 2.}
\end{omfigure}

There are two fixes. \code{BEGIN IMMEDIATE} takes the right to write at the start of the \omterm{def:b2:database-design:transaction}{transaction}, so the second desk
is refused until the first \omterm{def:b2:database-design:transaction}{commits}, and then reads the new value. Or the read and the write become one statement,
\code{UPDATE shelf SET available = available - 1}, which the database executes as a whole
(\cref{exo:b2:database-design:6}).

\begin{inpractice}
% ledger: b2:sqlite-begin-immediate
\code{BEGIN IMMEDIATE} starts a write \omterm{def:b2:database-design:transaction}{transaction} at once, and may fail with ``database is locked'' (\code{SQLITE_BUSY})
while another connection has one; a program then waits and tries again. A plain \code{BEGIN} waits for the first write,
and may fail there instead, after reading.
\end{inpractice}

Since SQLite never shows the other anomalies, a small simulator shows them: it runs a \omterm{def:b2:database-design:schedule}{schedule} of reads and writes under
three rules for answering a read, the latest write even if uncommitted, the latest committed value, or a snapshot of the
committed values when the \omterm{def:b2:database-design:transaction}{transaction} began. Under snapshots, a \omterm{def:b2:database-design:transaction}{transaction} that wants to \omterm{def:b2:database-design:transaction}{commit} a row that another
\omterm{def:b2:database-design:transaction}{transaction} committed since it began is refused, the rule of ``first committer wins''.

\omcode{code/bachelor-2/30-database-design/listings/schedules.py}{82}{102}{Two of the five staged \omterm{def:b2:database-design:schedule}{schedules}: the \omterm{def:b2:database-design:anomalies}{lost
update} and the \omterm{def:b2:database-design:anomalies}{write skew}}\label{lst:b2:database-design:schedules}

\omoutput{code/bachelor-2/30-database-design/outputs/schedules.txt}{Which anomaly each rule lets through.}\label{out:b2:database-design:schedules}

Reading committed values removes \omterm{def:b2:database-design:anomalies}{dirty reads} only; snapshots remove the three anomalies of the standard's table and turn
the \omterm{def:b2:database-design:anomalies}{lost update} into a refusal. \omterm{def:b2:database-design:anomalies}{Write skew} survives snapshots: the two \omterm{def:b2:database-design:transaction}{transactions} write different rows, so no write
conflicts with another, yet the rule ``at least one copy stays'' is broken. Only a serializable level prevents it, for
example by detecting that each read what the other wrote.

\section{Using an index}\label{sec:b2:database-design:index}

\begin{definition}[Index]\label{def:b2:database-design:index}
An \emph{index}\index{index (database)} on columns of a table is a separate structure, kept up to date by the database
system on every change, that maps values of those columns to the rows that hold them, in sorted order, so that the rows
with a given value or in a given range are found without reading the whole table.
\end{definition}

The database chooses for each query whether to use an \omterm{def:b2:database-design:index}{index}. \code{EXPLAIN QUERY PLAN} shows the choice:

\omcode{code/bachelor-2/30-database-design/listings/explain.sql}{1}{8}{The plan of the same query before and after an
\omterm{def:b2:database-design:index}{index}}\label{lst:b2:database-design:explain}

\omoutput{code/bachelor-2/30-database-design/outputs/explain.txt}{SCAN reads every row; SEARCH reads only the matching
ones, through the \omterm{def:b2:database-design:index}{index}.}\label{out:b2:database-design:explain}

The last query hides the column inside an expression, and the \omterm{def:b2:database-design:index}{index} no longer applies: \code{member_id + 0} is not a
value that the \omterm{def:b2:database-design:index}{index} sorts.

\begin{inpractice}
% ledger: b2:sqlite-eqp
SQLite's documentation says that SCAN is a full-table scan and SEARCH visits a subset of the rows, and warns that the
format of \code{EXPLAIN QUERY PLAN} may change between releases: it is for people to read, not for programs to parse.
\end{inpractice}

On 400 loans either way is instant. On generated tables of growing size:

\omcode{code/bachelor-2/30-database-design/listings/index_timing.py}{8}{34}{The loans of a member, by scan and by \omterm{def:b2:database-design:index}{index},
timed on generated tables}\label{lst:b2:database-design:index-timing}

\omoutput{code/bachelor-2/30-database-design/outputs/index_timing.txt}{Mean time of one query on the author's laptop,
in microseconds.}\label{out:b2:database-design:index-timing}

\begin{omfigure}
\begin{tikzpicture}
\begin{loglogaxis}[width=0.66\linewidth, height=5.2cm, xlabel={rows in the table}, ylabel={time of one query
  ($\mu$s)}, legend pos=north west, legend style={font=\small}, xtick={10000,100000,1000000}]
\addplot[mark=*, omThm] table[x=rows, y=scan] {code/bachelor-2/30-database-design/figures/index_timing.dat};
\addplot[mark=square*, omDef] table[x=rows, y=index] {code/bachelor-2/30-database-design/figures/index_timing.dat};
\legend{full scan, through the index}
\end{loglogaxis}
\end{tikzpicture}
\par\smallskip{\small \Cref{out:b2:database-design:index-timing} on logarithmic axes, measured on a laptop. The scan
grows in proportion to the table; the search through the \omterm{def:b2:database-design:index}{index} grows slowly, with the height of a tree, and with the
cost of reaching rows scattered in memory.}
\end{omfigure}

The scan time grows tenfold with each tenfold table; the \omterm{def:b2:database-design:index}{index} time grows about fivefold over a table a hundred times larger, and on a million rows the \omterm{def:b2:database-design:index}{index} is
about a thousand times faster. It is not free: building it took close to a second, and each later insertion or update of
a loan must update it too. How the \omterm{def:b2:database-design:index}{index} is organised, a balanced tree with many keys per node, and how the system
chooses among plans, are the subject of the Year 3 volume. An \omterm{def:b2:database-design:index}{index} does not always help either: a query that selects
a large fraction of the table reads it faster in order (\cref{exo:b2:database-design:12}).

\section{Exercises}\label{sec:b2:database-design:exercises}

\begin{exercise}[$\star$]\label{exo:b2:database-design:1}
With $F = \{a \to b,\ bc \to d,\ d \to e,\ ea \to c\}$, compute by hand the closures of $\{a\}$, $\{b, c\}$ and
$\{c, e\}$, then check with \code{exo_closure.py}. Is $\{a, d\}$ a superkey of $\{a, b, c, d, e\}$?
\end{exercise}

\begin{exercise}[$\star$]\label{exo:b2:database-design:2}
A BCNF step on the wide table's violation \emph{book} $\to$ \emph{title author genre pages} splits off only
\{book, title\}. What is wrong with the other part? Compare with the step that splits off the whole closure, using
\code{exo_bcnf_step.py}.
\end{exercise}

\begin{exercise}[$\star$]\label{exo:b2:database-design:3}
The reading room's bookings have attributes \emph{seat}, \emph{day}, \emph{slot}, \emph{member} and \emph{desk}, with
\emph{seat day slot} $\to$ \emph{member}, \emph{member day slot} $\to$ \emph{seat} and \emph{seat} $\to$ \emph{desk}. Find
all the candidate keys by hand, then with \code{exo_keys.py}.
\end{exercise}

\begin{exercise}[$\star$]\label{exo:b2:database-design:4}
A schema with $n$ attributes has exactly one candidate key, of size $k$. Prove that it has exactly $2^{n-k}$ superkeys.
\end{exercise}

\begin{exercise}[$\star\star$]\label{exo:b2:database-design:5}
Split the wide table into the member's attributes and the rest, keeping in both, in turn, the member, the street, or the
last name. Which splits are \omterm{def:b2:database-design:decomposition}{lossless} by Heath's theorem? For the others, explain how a join could invent rows. Check
with \code{exo_heath.py}, which runs the \emph{chase}: one row per part, with the symbol $a$ on the part's attributes and
distinct symbols elsewhere; the dependencies force equal symbols where two rows agree on a left side; the
\omterm{def:b2:database-design:decomposition}{decomposition} is \omterm{def:b2:database-design:decomposition}{lossless} exactly when some row ends with $a$ everywhere.
\end{exercise}

\begin{exercise}[$\star\star$]\label{exo:b2:database-design:6}
In the first part of \cref{lst:b2:database-design:lost-update}, explain the bug as a \omterm{def:b2:database-design:schedule}{schedule}, and say why each of the
two fixes prevents it. Which fix still works if the program must check that a copy is available before lending it?
\end{exercise}

\begin{exercise}[$\star\star$]\label{exo:b2:database-design:7}
Add a \emph{name} determined by the member to the bookings of \cref{exo:b2:database-design:3}, and let \emph{seat day
slot} determine \emph{desk} as well. Compute a \omterm{def:b2:database-design:cover}{minimal cover} and the 3NF synthesis by hand, then check with
\code{exo_synthesis.py}.
\end{exercise}

\begin{exercise}[$\star\star$]\label{exo:b2:database-design:8}
Derive the union rule and the \omterm{def:b2:database-design:decomposition}{decomposition} rule from Armstrong's three axioms.
\end{exercise}

\begin{exercise}[$\star\star$]\label{exo:b2:database-design:9}
Find, with \code{exo_extra_fds.py}, the dependencies between two single attributes that the town's 400 lines satisfy
but the library's rules do not imply. Which of them would you add to the schema, and what would change?
\end{exercise}

\begin{exercise}[$\star\star\star$]\label{exo:b2:database-design:10}
Prove \cref{thm:b2:database-design:heath}.
\end{exercise}

\begin{exercise}[$\star\star\star$]\label{exo:b2:database-design:11}
Prove \cref{prop:b2:database-design:closure-algo}.
\end{exercise}

\begin{exercise}[$\star\star\star$]\label{exo:b2:database-design:12}
A query selects the loans with \code{out_day} below a bound, through an \omterm{def:b2:database-design:index}{index} on \code{out_day} or by a full scan.
Measure both as the fraction selected grows, with \code{exo_selectivity.py}, find where the \omterm{def:b2:database-design:index}{index} stops helping, and
explain why.
\end{exercise}

\section{Problem: Redesigning the library's spreadsheet}\label{sec:b2:database-design:problem}

\begin{problem}[{Weekend problem --- Redesigning the library's spreadsheet}]\label{pb:b2:database-design:1}
The library replaces its spreadsheet by a database. The program below measures the old design, checks the new one and
replays member 55's move in both; \cref{out:b2:database-design:index-timing} measures the \omterm{def:b2:database-design:index}{index} the desks will need.

\omcode{code/bachelor-2/30-database-design/listings/redesign.py}{13}{39}{The old design measured, the new one
checked}\label{lst:b2:database-design:problem}

\omoutput{code/bachelor-2/30-database-design/outputs/redesign.txt}{Redundancy, violations, the four tables, and the move
in each design.}\label{out:b2:database-design:problem}

\needspace{4\baselineskip}\medskip\noindent\textbf{Part I --- The dependencies of the wide table.}\par\nopagebreak
\begin{enumerate}
\item List the library's rules as functional dependencies, and say which are about loans, books, members and streets.
\item Explain the 320 copies too many of the books' facts: what is stored 400 times, and what should be stored 80 times?
\item Give an \omterm{def:b2:database-design:anomaly}{update anomaly}, an \omterm{def:b2:database-design:anomaly}{insertion anomaly} and a \omterm{def:b2:database-design:anomaly}{deletion anomaly} of the wide table, with the library's own
facts.
\item Why is \emph{title} $\to$ \emph{book}, which the data satisfy, not one of the rules?
\item Does \emph{out} $\to$ \emph{due} hold? Should it be a rule, and what would that mean for the schema?
\item Compute the closure of \{member\} and of \{street\}.
\end{enumerate}

\needspace{4\baselineskip}\medskip\noindent\textbf{Part II --- Keys and normal forms.}\par\nopagebreak
\begin{enumerate}[resume, beginpenalty=10000]
\item Prove with the closure that \{loan\} is a candidate key of the wide table, and that it is the only one.
\item Which dependencies violate BCNF? Why exactly three left sides?
\item Is the wide table in 3NF? Use the definition.
\item Is the table (\emph{member}, \emph{first}, \emph{last}, \emph{street}) in BCNF? And (\emph{street},
\emph{postcode})?
\item What would a table (\emph{member}, \emph{street}, \emph{postcode}) violate?
\item Give the number of violations before and after the redesign, as the program prints them.
\end{enumerate}

\needspace{4\baselineskip}\medskip\noindent\textbf{Part III --- A lossless decomposition.}\par\nopagebreak
\begin{enumerate}[resume, beginpenalty=10000]
\item Draw the splits that lead to the four tables, naming the dependency used at each.
\item Prove that each split is \omterm{def:b2:database-design:decomposition}{lossless} with Heath's theorem.
\item Is the \omterm{def:b2:database-design:decomposition}{decomposition} dependency-preserving? Check each rule.
\item What did \cref{out:b2:database-design:decomposed} check on the town's data, and why are both counts needed?
\item Why does the streets table have 23 rows and not 24?
\item State the problem's result for the design.
\end{enumerate}

\needspace{4\baselineskip}\medskip\noindent\textbf{Part IV --- Transactions and an index.}\par\nopagebreak
\begin{enumerate}[resume, beginpenalty=10000]
\item How many rows does the move change in each design, and why can the new design not forget one?
\item A loan and its fine are recorded together. Why must they be one \omterm{def:b2:database-design:transaction}{transaction}?
\item Two desks lend the last copy of a book at the same moment. Which anomaly threatens, and how do you prevent it in
SQLite?
\item Which \omterm{def:b2:database-design:index}{index} do the desks need to list a member's loans, and what does \code{EXPLAIN QUERY PLAN} show before and
after it?
\item From \cref{out:b2:database-design:index-timing}, how much faster is the \omterm{def:b2:database-design:index}{index} on a million loans, and what does it
cost?
\item Complete the problem's result: from one wide table with three violating dependencies to four BCNF tables with
none, and an \omterm{def:b2:database-design:index}{index} that is how many times faster on a million loans?
\end{enumerate}
\end{problem}
